%%%%%%%%%%%%%%%%%%%%%%%%%%%%%%%%%%%%%%%%%%%%%%%%%%%%%%%%%%%%%%%%%%%%%%%%%%%%%%%%
% CHAPTER 13
%%%%%%%%%%%%%%%%%%%%%%%%%%%%%%%%%%%%%%%%%%%%%%%%%%%%%%%%%%%%%%%%%%%%%%%%%%%%%%%%

\chapter{Sampling Theorem and Discrete-Time Processing of Continuous-Time Signals}

\begin{tcolorbox}[colback=blue!5!white,colframe=blue!70!black,title=Learning Objectives]

After studying this chapter, the reader will be able to

\begin{itemize}
\item Understand the need for sampling.
\item Define sampling frequency and sampling period.
\item State and prove the Nyquist sampling theorem.
\item Understand aliasing and spectral overlap.
\item Determine the minimum sampling frequency.
\item Analyze ideal sampling in the time and frequency domains.
\end{itemize}

\end{tcolorbox}

%%%%%%%%%%%%%%%%%%%%%%%%%%%%%%%%%%%%%%%%%%%%%%%%%%%%%%%%%%%%%%%%%%%%%%%%%%%%%%%%

\section{Introduction}

Real-world signals such as speech, ECG, radar and sensor outputs are
\textbf{continuous-time signals}.

Digital systems cannot process continuous-time signals directly.

Therefore, we must convert them into discrete-time signals by
\textbf{sampling}.

Sampling is the bridge between the analog world and the digital world.

%%%%%%%%%%%%%%%%%%%%%%%%%%%%%%%%%%%%%%%%%%%%%%%%%%%%%%%%%%%%%%%%%%%%%%%%%%%%%%%%

\section{Continuous-Time Signal}

A continuous-time signal is defined for every real value of time:

\[
x(t),\qquad -\infty<t<\infty.
\]

Example:

\[
x(t)=A\cos(2\pi f_0 t).
\]

%%%%%%%%%%%%%%%%%%%%%%%%%%%%%%%%%%%%%%%%%%%%%%%%%%%%%%%%%%%%%%%%%%%%%%%%%%%%%%%%

\section{Sampling Process}

Sampling means measuring the signal at equally spaced instants:

\[
t=nT_s,
\qquad n=0,\pm1,\pm2,\ldots
\]

The sampled sequence is

\[
\boxed{
x[n]=x(nT_s)
}
\]

where

\begin{itemize}
\item \(T_s\): sampling period,
\item \(f_s=\dfrac{1}{T_s}\): sampling frequency.
\end{itemize}

%%%%%%%%%%%%%%%%%%%%%%%%%%%%%%%%%%%%%%%%%%%%%%%%%%%%%%%%%%%%%%%%%%%%%%%%%%%%%%%%

\section{Example}

If

\[
x(t)=\cos(200\pi t)
\]

and

\[
f_s=1000\text{ Hz},
\]

then

\[
T_s=1\text{ ms}.
\]

The samples are

\[
x[n]=\cos(200\pi nT_s)
=\cos(0.2\pi n).
\]

%%%%%%%%%%%%%%%%%%%%%%%%%%%%%%%%%%%%%%%%%%%%%%%%%%%%%%%%%%%%%%%%%%%%%%%%%%%%%%%%

\section{Impulse-Train Representation}

Ideal sampling is modeled by multiplying \(x(t)\) with an impulse train:

\[
p(t)=
\sum_{n=-\infty}^{\infty}\delta(t-nT_s).
\]

The sampled signal is

\[
\boxed{
x_s(t)=x(t)p(t)
=
\sum_{n=-\infty}^{\infty}
x(nT_s)\delta(t-nT_s)
}
\]

%%%%%%%%%%%%%%%%%%%%%%%%%%%%%%%%%%%%%%%%%%%%%%%%%%%%%%%%%%%%%%%%%%%%%%%%%%%%%%%%

\section{Spectrum of the Impulse Train}

The Fourier transform of \(p(t)\) is another impulse train:

\[
\boxed{
P(f)=
f_s
\sum_{k=-\infty}^{\infty}
\delta(f-kf_s)
}
\]

Thus, the sampling operation creates periodic replicas of the original spectrum.

%%%%%%%%%%%%%%%%%%%%%%%%%%%%%%%%%%%%%%%%%%%%%%%%%%%%%%%%%%%%%%%%%%%%%%%%%%%%%%%%

\section{Spectrum of the Sampled Signal}

Since multiplication in time corresponds to convolution in frequency,

\[
X_s(f)=X(f)*P(f).
\]

Substituting \(P(f)\),

\[
\boxed{
X_s(f)=
f_s
\sum_{k=-\infty}^{\infty}
X(f-kf_s)
}
\]

This means the original spectrum is repeated every \(f_s\) hertz.

%%%%%%%%%%%%%%%%%%%%%%%%%%%%%%%%%%%%%%%%%%%%%%%%%%%%%%%%%%%%%%%%%%%%%%%%%%%%%%%%

\section{Bandlimited Signal}

A signal is bandlimited if

\[
X(f)=0
\qquad \text{for } |f|>B.
\]

Here,

\[
B
\]

is the highest frequency component.

%%%%%%%%%%%%%%%%%%%%%%%%%%%%%%%%%%%%%%%%%%%%%%%%%%%%%%%%%%%%%%%%%%%%%%%%%%%%%%%%

\section{Nyquist Sampling Theorem}

\begin{tcolorbox}[colback=green!5!white,colframe=green!60!black,title=Nyquist Theorem]

A bandlimited signal with highest frequency \(B\) can be reconstructed exactly
from its samples if

\[
\boxed{
f_s\ge 2B
}
\]

\end{tcolorbox}

The quantity

\[
\boxed{
2B
}
\]

is called the \textbf{Nyquist rate}.

%%%%%%%%%%%%%%%%%%%%%%%%%%%%%%%%%%%%%%%%%%%%%%%%%%%%%%%%%%%%%%%%%%%%%%%%%%%%%%%%

\section{Why \(2B\)?}

The spectral replicas are centered at

\[
0,\pm f_s,\pm2f_s,\ldots
\]

To avoid overlap,

\[
f_s-B\ge B.
\]

Hence,

\[
f_s\ge2B.
\]

%%%%%%%%%%%%%%%%%%%%%%%%%%%%%%%%%%%%%%%%%%%%%%%%%%%%%%%%%%%%%%%%%%%%%%%%%%%%%%%%

\section{Aliasing}

If

\[
f_s<2B,
\]

the spectral replicas overlap.

This overlap is called

\[
\boxed{
\text{Aliasing}
}
\]

Aliasing causes high-frequency components to appear as lower frequencies after
sampling.

%%%%%%%%%%%%%%%%%%%%%%%%%%%%%%%%%%%%%%%%%%%%%%%%%%%%%%%%%%%%%%%%%%%%%%%%%%%%%%%%

\section{Aliasing Formula}

A frequency \(f\) appears after sampling as

\[
\boxed{
f_a=
|f-kf_s|
}
\]

for some integer \(k\) chosen so that

\[
0\le f_a\le \frac{f_s}{2}.
\]

%%%%%%%%%%%%%%%%%%%%%%%%%%%%%%%%%%%%%%%%%%%%%%%%%%%%%%%%%%%%%%%%%%%%%%%%%%%%%%%%

\section{Example 1}

A 900 Hz sinusoid is sampled at

\[
f_s=1000\text{ Hz}.
\]

%%%%%%%%%%%%%%%%%%%%%%%%%%%%%%%%%%%%%%%%%%%%%%%%%%%%%%%%%%%%%%%%%%%%%%%%%%%%%%%%

\subsection{Compute Aliased Frequency}

Choose \(k=1\):

\[
f_a=|900-1000|=100\text{ Hz}.
\]

\[
\boxed{
900\text{ Hz aliases to }100\text{ Hz}
}
\]

%%%%%%%%%%%%%%%%%%%%%%%%%%%%%%%%%%%%%%%%%%%%%%%%%%%%%%%%%%%%%%%%%%%%%%%%%%%%%%%%

\section{Nyquist Frequency}

The highest frequency that can be represented without ambiguity is

\[
\boxed{
f_N=\frac{f_s}{2}
}
\]

called the \textbf{Nyquist frequency}.

%%%%%%%%%%%%%%%%%%%%%%%%%%%%%%%%%%%%%%%%%%%%%%%%%%%%%%%%%%%%%%%%%%%%%%%%%%%%%%%%

\section{Example 2}

If

\[
f_s=8\text{ kHz},
\]

then

\[
\boxed{
f_N=4\text{ kHz}
}
\]

Any component above 4 kHz will alias.

%%%%%%%%%%%%%%%%%%%%%%%%%%%%%%%%%%%%%%%%%%%%%%%%%%%%%%%%%%%%%%%%%%%%%%%%%%%%%%%%

\section{Anti-Aliasing Filter}

Before sampling, an analog low-pass filter is used.

\begin{center}
\begin{tikzpicture}[node distance=2.2cm,>=Latex]
\node[draw,minimum width=2.4cm,minimum height=0.9cm] (x) {Input \(x(t)\)};
\node[draw,minimum width=3.2cm,minimum height=0.9cm,right of=x] (lpf) {Anti-aliasing LPF};
\node[draw,minimum width=2.6cm,minimum height=0.9cm,right of=lpf] (sampler) {Sampler};

\draw[->] (x) -- (lpf);
\draw[->] (lpf) -- (sampler);
\end{tikzpicture}
\end{center}

The filter removes frequency components above \(B\).

%%%%%%%%%%%%%%%%%%%%%%%%%%%%%%%%%%%%%%%%%%%%%%%%%%%%%%%%%%%%%%%%%%%%%%%%%%%%%%%%

\section{Ideal Reconstruction}

The original signal can be reconstructed using an ideal low-pass filter with
cutoff \(B\).

%%%%%%%%%%%%%%%%%%%%%%%%%%%%%%%%%%%%%%%%%%%%%%%%%%%%%%%%%%%%%%%%%%%%%%%%%%%%%%%%

\subsection{Interpolation Formula}

\[
\boxed{
x(t)=
\sum_{n=-\infty}^{\infty}
x[n]
\operatorname{sinc}
\left(
\frac{t-nT_s}{T_s}
\right)
}
\]

where

\[
\operatorname{sinc}(x)=
\frac{\sin(\pi x)}{\pi x}.
\]

%%%%%%%%%%%%%%%%%%%%%%%%%%%%%%%%%%%%%%%%%%%%%%%%%%%%%%%%%%%%%%%%%%%%%%%%%%%%%%%%

\section{Interpretation of Sinc Interpolation}

Each sample generates a sinc pulse.

\begin{itemize}
\item The pulse equals 1 at its own sampling instant.
\item It equals 0 at all other sampling instants.
\item Adding all sinc pulses reconstructs the original signal.
\end{itemize}

%%%%%%%%%%%%%%%%%%%%%%%%%%%%%%%%%%%%%%%%%%%%%%%%%%%%%%%%%%%%%%%%%%%%%%%%%%%%%%%%

\section{Under-Sampling}

If

\[
f_s<2B,
\]

the signal is under-sampled.

Example:

\[
B=5\text{ kHz},
\qquad
f_s=8\text{ kHz}.
\]

Since

\[
8<10,
\]

the signal cannot be reconstructed exactly.

%%%%%%%%%%%%%%%%%%%%%%%%%%%%%%%%%%%%%%%%%%%%%%%%%%%%%%%%%%%%%%%%%%%%%%%%%%%%%%%%

\section{Critical Sampling}

When

\[
f_s=2B,
\]

the spectral replicas just touch each other.

This is called \textbf{critical sampling}.

In practice, we usually choose

\[
f_s>2B
\]

to allow a transition band for the anti-aliasing filter.

%%%%%%%%%%%%%%%%%%%%%%%%%%%%%%%%%%%%%%%%%%%%%%%%%%%%%%%%%%%%%%%%%%%%%%%%%%%%%%%%

\section{Over-Sampling}

When

\[
f_s\gg2B,
\]

the system is over-sampled.

Advantages:

\begin{itemize}
\item easier analog filter design,
\item improved noise performance,
\item higher dynamic range.
\end{itemize}

%%%%%%%%%%%%%%%%%%%%%%%%%%%%%%%%%%%%%%%%%%%%%%%%%%%%%%%%%%%%%%%%%%%%%%%%%%%%%%%%

\section{Solved Example 13.1}

A signal is bandlimited to

\[
B=3\text{ kHz}.
\]

Find the minimum sampling frequency.

%%%%%%%%%%%%%%%%%%%%%%%%%%%%%%%%%%%%%%%%%%%%%%%%%%%%%%%%%%%%%%%%%%%%%%%%%%%%%%%%

\subsection{Solution}

\[
f_s^{\min}=2B=6\text{ kHz}.
\]

\[
\boxed{
f_s^{\min}=6\text{ kHz}
}
\]

%%%%%%%%%%%%%%%%%%%%%%%%%%%%%%%%%%%%%%%%%%%%%%%%%%%%%%%%%%%%%%%%%%%%%%%%%%%%%%%%

\section{Solved Example 13.2}

A 7 kHz signal is sampled at

\[
f_s=10\text{ kHz}.
\]

Find the aliased frequency.

%%%%%%%%%%%%%%%%%%%%%%%%%%%%%%%%%%%%%%%%%%%%%%%%%%%%%%%%%%%%%%%%%%%%%%%%%%%%%%%%

\subsection{Solution}

\[
f_a=|7-10|=3\text{ kHz}.
\]

\[
\boxed{
f_a=3\text{ kHz}
}
\]

%%%%%%%%%%%%%%%%%%%%%%%%%%%%%%%%%%%%%%%%%%%%%%%%%%%%%%%%%%%%%%%%%%%%%%%%%%%%%%%%

\section{Solved Example 13.3}

A signal contains frequencies up to

\[
15\text{ kHz}.
\]

Choose a practical sampling frequency.

%%%%%%%%%%%%%%%%%%%%%%%%%%%%%%%%%%%%%%%%%%%%%%%%%%%%%%%%%%%%%%%%%%%%%%%%%%%%%%%%

\subsection{Solution}

Nyquist rate:

\[
2B=30\text{ kHz}.
\]

Choose a slightly higher value:

\[
\boxed{
f_s=36\text{ kHz}
}
\]

(or 40 kHz in practice).

%%%%%%%%%%%%%%%%%%%%%%%%%%%%%%%%%%%%%%%%%%%%%%%%%%%%%%%%%%%%%%%%%%%%%%%%%%%%%%%%

\section{Solved Example 13.4}

Find the discrete-time frequency corresponding to

\[
f=1\text{ kHz},
\qquad
f_s=8\text{ kHz}.
\]

%%%%%%%%%%%%%%%%%%%%%%%%%%%%%%%%%%%%%%%%%%%%%%%%%%%%%%%%%%%%%%%%%%%%%%%%%%%%%%%%

\subsection{Solution}

\[
\omega=
2\pi\frac{f}{f_s}
=
2\pi\frac{1000}{8000}
=
\frac{\pi}{4}.
\]

\[
\boxed{
\omega=\frac{\pi}{4}\text{ rad/sample}
}
\]

%%%%%%%%%%%%%%%%%%%%%%%%%%%%%%%%%%%%%%%%%%%%%%%%%%%%%%%%%%%%%%%%%%%%%%%%%%%%%%%%

\section{Quick Check}

\begin{enumerate}
\item \(T_s=\boxed{1/f_s}\).
\item Nyquist rate \(=\boxed{2B}\).
\item Nyquist frequency \(=\boxed{f_s/2}\).
\item Aliasing occurs when \(\boxed{f_s<2B}\).
\item Reconstruction uses \(\boxed{\text{sinc interpolation}}\).
\end{enumerate}

%%%%%%%%%%%%%%%%%%%%%%%%%%%%%%%%%%%%%%%%%%%%%%%%%%%%%%%%%%%%%%%%%%%%%%%%%%%%%%%%

\begin{tcolorbox}[title=One-Minute Revision,colback=blue!5]

\[
x[n]=x(nT_s)
\]

\[
f_s=\frac{1}{T_s}
\]

\[
x_s(t)=
\sum x(nT_s)\delta(t-nT_s)
\]

\[
X_s(f)=
f_s\sum X(f-kf_s)
\]

Nyquist theorem:

\[
f_s\ge2B
\]

Nyquist frequency:

\[
f_N=\frac{f_s}{2}
\]

Aliasing:

\[
f_a=|f-kf_s|
\]

Reconstruction:

\[
x(t)=
\sum x[n]\,
\operatorname{sinc}
\left(
\frac{t-nT_s}{T_s}
\right)
\]

\end{tcolorbox}
%%%%%%%%%%%%%%%%%%%%%%%%%%%%%%%%%%%%%%%%%%%%%%%%%%%%%%%%%%%%%%%%%%%%%%%%%%%%%%%%
%%%%%%%%%%%%%%%%%%%%%%%%%%%%%%%%%%%%%%%%%%%%%%%%%%%%%%%%%%%%%%%%%%%%%%%%%%%%%%%%
\section{Discrete-Time Processing of Continuous-Time Signals}

A practical DSP system consists of three stages:

\begin{center}
\begin{tikzpicture}[node distance=2.2cm,>=Latex]
\node[draw,minimum width=2.2cm,minimum height=0.9cm] (x) {Analog \(x(t)\)};
\node[draw,minimum width=1.8cm,minimum height=0.9cm,right of=x] (adc) {ADC};
\node[draw,minimum width=2.2cm,minimum height=0.9cm,right of=adc] (dsp) {DSP};
\node[draw,minimum width=1.8cm,minimum height=0.9cm,right of=dsp] (dac) {DAC};
\node[draw,minimum width=2.4cm,minimum height=0.9cm,right of=dac] (y) {Analog \(y(t)\)};

\draw[->] (x) -- (adc);
\draw[->] (adc) -- (dsp);
\draw[->] (dsp) -- (dac);
\draw[->] (dac) -- (y);
\end{tikzpicture}
\end{center}

\begin{itemize}
\item \textbf{ADC}: sampling and quantization.
\item \textbf{DSP}: digital filtering, modulation, FFT, etc.
\item \textbf{DAC}: converts the processed sequence back to analog form.
\end{itemize}

%%%%%%%%%%%%%%%%%%%%%%%%%%%%%%%%%%%%%%%%%%%%%%%%%%%%%%%%%%%%%%%%%%%%%%%%%%%%%%%%

\section{Analog-to-Digital Conversion (ADC)}

The ADC performs two operations:

%%%%%%%%%%%%%%%%%%%%%%%%%%%%%%%%%%%%%%%%%%%%%%%%%%%%%%%%%%%%%%%%%%%%%%%%%%%%%%%%

\subsection{Sampling}

\[
x[n]=x(nT_s)
\]

%%%%%%%%%%%%%%%%%%%%%%%%%%%%%%%%%%%%%%%%%%%%%%%%%%%%%%%%%%%%%%%%%%%%%%%%%%%%%%%%

\subsection{Quantization}

Each sample is rounded to one of a finite number of amplitude levels.

%%%%%%%%%%%%%%%%%%%%%%%%%%%%%%%%%%%%%%%%%%%%%%%%%%%%%%%%%%%%%%%%%%%%%%%%%%%%%%%%

\section{Quantization Levels}

For an \(N_b\)-bit ADC,

\[
\boxed{
L=2^{N_b}
}
\]

levels are available.

%%%%%%%%%%%%%%%%%%%%%%%%%%%%%%%%%%%%%%%%%%%%%%%%%%%%%%%%%%%%%%%%%%%%%%%%%%%%%%%%

\subsection{Examples}

\begin{table}[H]
\centering
\caption{ADC resolution}
\begin{tabular}{|c|c|}
\hline
\textbf{Bits} & \textbf{Levels}\\
\hline
8 & 256\\
10 & 1024\\
12 & 4096\\
16 & 65536\\
\hline
\end{tabular}
\end{table}

%%%%%%%%%%%%%%%%%%%%%%%%%%%%%%%%%%%%%%%%%%%%%%%%%%%%%%%%%%%%%%%%%%%%%%%%%%%%%%%%

\section{Quantization Step Size}

If the input range is

\[
[-V_{\max},\,V_{\max}],
\]

the step size is

\[
\boxed{
\Delta=
\frac{2V_{\max}}{2^{N_b}}
}
\]

%%%%%%%%%%%%%%%%%%%%%%%%%%%%%%%%%%%%%%%%%%%%%%%%%%%%%%%%%%%%%%%%%%%%%%%%%%%%%%%%

\section{Quantization Error}

The error is

\[
e_q=x-x_q.
\]

For an ideal uniform quantizer,

\[
\boxed{
-\frac{\Delta}{2}\le e_q\le\frac{\Delta}{2}
}
\]

and the mean-square error is

\[
\boxed{
P_q=\frac{\Delta^2}{12}.
}
\]

%%%%%%%%%%%%%%%%%%%%%%%%%%%%%%%%%%%%%%%%%%%%%%%%%%%%%%%%%%%%%%%%%%%%%%%%%%%%%%%%

\section{Signal-to-Quantization-Noise Ratio}

For a full-scale sinusoid,

\[
\boxed{
\text{SQNR}_{\text{dB}}
\approx
6.02N_b+1.76
}
\]

%%%%%%%%%%%%%%%%%%%%%%%%%%%%%%%%%%%%%%%%%%%%%%%%%%%%%%%%%%%%%%%%%%%%%%%%%%%%%%%%

\subsection{Example}

For a 12-bit ADC,

\[
\text{SQNR}\approx
6.02(12)+1.76
=
74\text{ dB}.
\]

%%%%%%%%%%%%%%%%%%%%%%%%%%%%%%%%%%%%%%%%%%%%%%%%%%%%%%%%%%%%%%%%%%%%%%%%%%%%%%%%

\section{Digital-to-Analog Conversion (DAC)}

The DAC produces a staircase waveform.

An analog low-pass filter then smooths the staircase to reconstruct the continuous
signal.

%%%%%%%%%%%%%%%%%%%%%%%%%%%%%%%%%%%%%%%%%%%%%%%%%%%%%%%%%%%%%%%%%%%%%%%%%%%%%%%%

\section{Zero-Order Hold (ZOH)}

A practical DAC holds each sample value constant for one sampling period.

The ZOH output is

\[
x_{\text{ZOH}}(t)=x[n],
\qquad nT_s\le t<(n+1)T_s.
\]

%%%%%%%%%%%%%%%%%%%%%%%%%%%%%%%%%%%%%%%%%%%%%%%%%%%%%%%%%%%%%%%%%%%%%%%%%%%%%%%%

\section{ZOH Frequency Response}

The ZOH introduces a sinc-shaped attenuation:

\[
\boxed{
H_{\text{ZOH}}(f)=
T_s\,\mathrm{sinc}(fT_s)\,
e^{-j\pi fT_s}
}
\]

High frequencies are attenuated more than low frequencies.

%%%%%%%%%%%%%%%%%%%%%%%%%%%%%%%%%%%%%%%%%%%%%%%%%%%%%%%%%%%%%%%%%%%%%%%%%%%%%%%%

\section{Sampling of a Sinusoid}

Let

\[
x(t)=A\cos(2\pi f_0 t+\phi).
\]

Sampling at \(f_s\) gives

\[
x[n]=A\cos\left(2\pi\frac{f_0}{f_s}n+\phi\right).
\]

Define the discrete-time frequency

\[
\boxed{
\omega_0=
2\pi\frac{f_0}{f_s}
}
\]

in rad/sample.

%%%%%%%%%%%%%%%%%%%%%%%%%%%%%%%%%%%%%%%%%%%%%%%%%%%%%%%%%%%%%%%%%%%%%%%%%%%%%%%%

\section{Example 1}

\[
f_0=1\text{ kHz},
\qquad
f_s=8\text{ kHz}.
\]

Then

\[
\omega_0=
2\pi\frac{1}{8}
=
\frac{\pi}{4}.
\]

Hence,

\[
\boxed{
x[n]=A\cos\left(\frac{\pi}{4}n+\phi\right)
}
\]

%%%%%%%%%%%%%%%%%%%%%%%%%%%%%%%%%%%%%%%%%%%%%%%%%%%%%%%%%%%%%%%%%%%%%%%%%%%%%%%%

\section{Frequency Equivalence}

Discrete-time frequencies separated by multiples of \(2\pi\) are identical:

\[
\boxed{
\omega\equiv\omega+2\pi k
}
\]

Therefore,

\[
\cos(\omega n)=\cos((\omega+2\pi)n).
\]

%%%%%%%%%%%%%%%%%%%%%%%%%%%%%%%%%%%%%%%%%%%%%%%%%%%%%%%%%%%%%%%%%%%%%%%%%%%%%%%%

\section{Example 2}

\[
\omega_1=\frac{\pi}{4},
\qquad
\omega_2=\frac{9\pi}{4}.
\]

Since

\[
\omega_2=\omega_1+2\pi,
\]

both sequences are identical.

%%%%%%%%%%%%%%%%%%%%%%%%%%%%%%%%%%%%%%%%%%%%%%%%%%%%%%%%%%%%%%%%%%%%%%%%%%%%%%%%

\section{Analog Frequency Mapping}

The analog frequency corresponding to a discrete-time frequency \(\omega\) is

\[
\boxed{
f=\frac{\omega}{2\pi}f_s
}
\]

%%%%%%%%%%%%%%%%%%%%%%%%%%%%%%%%%%%%%%%%%%%%%%%%%%%%%%%%%%%%%%%%%%%%%%%%%%%%%%%%

\subsection{Example}

\[
\omega=\frac{\pi}{2},
\qquad
f_s=10\text{ kHz}.
\]

Then

\[
f=
\frac{\pi/2}{2\pi}(10\,000)
=
2500\text{ Hz}.
\]

\[
\boxed{
f=2.5\text{ kHz}
}
\]

%%%%%%%%%%%%%%%%%%%%%%%%%%%%%%%%%%%%%%%%%%%%%%%%%%%%%%%%%%%%%%%%%%%%%%%%%%%%%%%%

\section{Discrete-Time Filtering of Analog Signals}

Suppose the analog input is sampled to obtain \(x[n]\).

A digital filter with frequency response \(H(e^{j\omega})\) produces

\[
y[n]=x[n]*h[n].
\]

After DAC and reconstruction, the analog output approximates the response of an
analog filter.

%%%%%%%%%%%%%%%%%%%%%%%%%%%%%%%%%%%%%%%%%%%%%%%%%%%%%%%%%%%%%%%%%%%%%%%%%%%%%%%%

\section{Equivalent Analog Frequency}

The digital filter acts on frequencies

\[
0\le\omega\le\pi.
\]

These correspond to analog frequencies

\[
0\le f\le\frac{f_s}{2}.
\]

Thus, digital filter specifications are often converted between \(\omega\) and
\(f\).

%%%%%%%%%%%%%%%%%%%%%%%%%%%%%%%%%%%%%%%%%%%%%%%%%%%%%%%%%%%%%%%%%%%%%%%%%%%%%%%%

\section{Multirate Signal Processing}

%%%%%%%%%%%%%%%%%%%%%%%%%%%%%%%%%%%%%%%%%%%%%%%%%%%%%%%%%%%%%%%%%%%%%%%%%%%%%%%%

\subsection{Decimation}

Reducing the sampling rate by \(M\):

\[
\boxed{
y[n]=x[nM]
}
\]

The new sampling frequency is

\[
f_s'=\frac{f_s}{M}.
\]

An anti-aliasing filter must be applied before decimation.

%%%%%%%%%%%%%%%%%%%%%%%%%%%%%%%%%%%%%%%%%%%%%%%%%%%%%%%%%%%%%%%%%%%%%%%%%%%%%%%%

\section{Interpolation}

Increasing the sampling rate by \(L\):

\begin{enumerate}
\item Insert \(L-1\) zeros between samples.
\item Apply a low-pass interpolation filter.
\end{enumerate}

The new sampling frequency is

\[
f_s'=Lf_s.
\]

%%%%%%%%%%%%%%%%%%%%%%%%%%%%%%%%%%%%%%%%%%%%%%%%%%%%%%%%%%%%%%%%%%%%%%%%%%%%%%%%

\section{Example 3: Decimation}

A signal sampled at

\[
48\text{ kHz}
\]

is decimated by

\[
M=4.
\]

Then

\[
f_s'=
\frac{48}{4}
=
12\text{ kHz}.
\]

\[
\boxed{
f_s'=12\text{ kHz}
}
\]

%%%%%%%%%%%%%%%%%%%%%%%%%%%%%%%%%%%%%%%%%%%%%%%%%%%%%%%%%%%%%%%%%%%%%%%%%%%%%%%%

\section{Example 4: Interpolation}

A signal sampled at

\[
8\text{ kHz}
\]

is interpolated by

\[
L=3.
\]

Then

\[
f_s'=24\text{ kHz}.
\]

\[
\boxed{
f_s'=24\text{ kHz}
}
\]

%%%%%%%%%%%%%%%%%%%%%%%%%%%%%%%%%%%%%%%%%%%%%%%%%%%%%%%%%%%%%%%%%%%%%%%%%%%%%%%%

\section{Bandpass Sampling (Concept)}

A bandpass signal can sometimes be sampled at a rate lower than \(2f_H\), provided
the spectral replicas do not overlap.

This is called \textbf{undersampling} or \textbf{bandpass sampling} and is used in
RF receivers.

For GATE, remember:

\[
\boxed{
2B\le f_s
}
\]

where \(B\) is the bandwidth of the bandpass signal.

%%%%%%%%%%%%%%%%%%%%%%%%%%%%%%%%%%%%%%%%%%%%%%%%%%%%%%%%%%%%%%%%%%%%%%%%%%%%%%%%

\section{Solved Example 13.5}

A signal is bandlimited to

\[
B=5\text{ kHz}.
\]

Find the minimum sampling frequency.

%%%%%%%%%%%%%%%%%%%%%%%%%%%%%%%%%%%%%%%%%%%%%%%%%%%%%%%%%%%%%%%%%%%%%%%%%%%%%%%%

\subsection{Solution}

\[
f_s^{\min}=2B=10\text{ kHz}.
\]

\[
\boxed{
10\text{ kHz}
}
\]

%%%%%%%%%%%%%%%%%%%%%%%%%%%%%%%%%%%%%%%%%%%%%%%%%%%%%%%%%%%%%%%%%%%%%%%%%%%%%%%%

\section{Solved Example 13.6}

A 9 kHz sinusoid is sampled at

\[
f_s=12\text{ kHz}.
\]

Find the aliased frequency.

%%%%%%%%%%%%%%%%%%%%%%%%%%%%%%%%%%%%%%%%%%%%%%%%%%%%%%%%%%%%%%%%%%%%%%%%%%%%%%%%

\subsection{Solution}

Choose \(k=1\):

\[
f_a=|9-12|=3\text{ kHz}.
\]

\[
\boxed{
f_a=3\text{ kHz}
}
\]

%%%%%%%%%%%%%%%%%%%%%%%%%%%%%%%%%%%%%%%%%%%%%%%%%%%%%%%%%%%%%%%%%%%%%%%%%%%%%%%%

\section{Solved Example 13.7}

Find the discrete-time frequency of a 2 kHz signal sampled at 16 kHz.

%%%%%%%%%%%%%%%%%%%%%%%%%%%%%%%%%%%%%%%%%%%%%%%%%%%%%%%%%%%%%%%%%%%%%%%%%%%%%%%%

\subsection{Solution}

\[
\omega=
2\pi\frac{2000}{16000}
=
\frac{\pi}{4}.
\]

\[
\boxed{
\omega=\frac{\pi}{4}\text{ rad/sample}
}
\]

%%%%%%%%%%%%%%%%%%%%%%%%%%%%%%%%%%%%%%%%%%%%%%%%%%%%%%%%%%%%%%%%%%%%%%%%%%%%%%%%

\section{Solved Example 13.8}

A 10-bit ADC has input range

\[
\pm5\text{ V}.
\]

Find the quantization step size.

%%%%%%%%%%%%%%%%%%%%%%%%%%%%%%%%%%%%%%%%%%%%%%%%%%%%%%%%%%%%%%%%%%%%%%%%%%%%%%%%

\subsection{Solution}

\[
\Delta=
\frac{10}{2^{10}}
=
\frac{10}{1024}
=
9.77\text{ mV}.
\]

\[
\boxed{
\Delta\approx9.77\text{ mV}
}
\]

%%%%%%%%%%%%%%%%%%%%%%%%%%%%%%%%%%%%%%%%%%%%%%%%%%%%%%%%%%%%%%%%%%%%%%%%%%%%%%%%

\section{Solved Example 13.9}

Find the SQNR of an 8-bit ADC.

%%%%%%%%%%%%%%%%%%%%%%%%%%%%%%%%%%%%%%%%%%%%%%%%%%%%%%%%%%%%%%%%%%%%%%%%%%%%%%%%

\subsection{Solution}

\[
\text{SQNR}\approx
6.02(8)+1.76
=
49.92\text{ dB}.
\]

\[
\boxed{
\text{SQNR}\approx50\text{ dB}
}
\]

%%%%%%%%%%%%%%%%%%%%%%%%%%%%%%%%%%%%%%%%%%%%%%%%%%%%%%%%%%%%%%%%%%%%%%%%%%%%%%%%

\section{Quick Check}

\begin{enumerate}
\item ADC levels \(=\boxed{2^{N_b}}\).
\item Quantization step \(=\boxed{\dfrac{2V_{\max}}{2^{N_b}}}\).
\item SQNR \(=\boxed{6.02N_b+1.76}\).
\item Discrete-time frequency:
\[
\boxed{\omega=2\pi f/f_s}
\]
\item Decimation:
\[
\boxed{y[n]=x[nM]}
\]
\item Interpolation increases sampling rate by \(\boxed{L}\).
\end{enumerate}

%%%%%%%%%%%%%%%%%%%%%%%%%%%%%%%%%%%%%%%%%%%%%%%%%%%%%%%%%%%%%%%%%%%%%%%%%%%%%%%%

\section{Chapter Summary}

In this chapter, we studied the complete chain of discrete-time processing of
continuous-time signals, including sampling, ADC, quantization, DAC, zero-order
hold, frequency mapping and multirate processing.

We also analyzed aliasing, reconstruction, decimation and interpolation, and
solved several GATE-style numerical problems.

With this chapter, the entire GATE Signals \& Systems syllabus is complete.

%%%%%%%%%%%%%%%%%%%%%%%%%%%%%%%%%%%%%%%%%%%%%%%%%%%%%%%%%%%%%%%%%%%%%%%%%%%%%%%%

\section{Complete Volume Conclusion}

This volume covered

\begin{enumerate}
\item Continuous-time Fourier series,
\item Continuous-time Fourier transform,
\item Sampling theorem,
\item DTFT,
\item DFT,
\item FFT,
\item z-transform,
\item Frequency response,
\item Poles and zeros,
\item Phase delay,
\item Group delay,
\item Discrete-time processing of continuous-time signals.
\end{enumerate}

These topics form the mathematical foundation of

\begin{itemize}
\item Digital Signal Processing,
\item Communication Systems,
\item Control Systems,
\item Embedded DSP,
\item Radar and Sonar,
\item Audio and Speech Processing,
\item Satellite and Aerospace Signal Processing.
\end{itemize}

%%%%%%%%%%%%%%%%%%%%%%%%%%%%%%%%%%%%%%%%%%%%%%%%%%%%%%%%%%%%%%%%%%%%%%%%%%%%%%%%

\begin{tcolorbox}[title=Final One-Minute Revision,colback=blue!5]

Sampling:

\[
x[n]=x(nT_s),
\qquad
f_s=\frac1{T_s}
\]

Nyquist:

\[
f_s\ge2B
\]

Aliasing:

\[
f_a=|f-kf_s|
\]

DTFT:

\[
X(e^{j\omega})=
\sum x[n]e^{-j\omega n}
\]

DFT:

\[
X[k]=
\sum x[n]e^{-j2\pi kn/N}
\]

FFT complexity:

\[
N\log_2N
\]

z-transform:

\[
X(z)=
\sum x[n]z^{-n}
\]

Frequency response:

\[
H(e^{j\omega})=H(z)|_{z=e^{j\omega}}
\]

Group delay:

\[
\tau_g=
-\frac{d\phi}{d\omega}
\]

Linear-phase FIR:

\[
\tau_g=\frac{N-1}{2}
\]

ADC:

\[
L=2^{N_b},
\qquad
\Delta=\frac{2V_{\max}}{2^{N_b}}
\]

SQNR:

\[
6.02N_b+1.76\text{ dB}
\]

\end{tcolorbox}
%%%%%%%%%%%%%%%%%%%%%%%%%%%%%%%%%%%%%%%%%%%%%%%%%%%%%%%%%%%%%%%%%%%%%%%%%%%%%%%%